%% 25장 행렬 함수와 행렬 방정식
\chapter{Matrix Functions and Matrix Equations}
\label{ch:matrix-functions}

\abstract{스칼라 함수 $f$를 행렬에 적용한 $f(A)$는 미분방정식, 마르코프 과정, 그래프의 분석에서 자연스럽게 나타나며, 그 가운데 가장 중요한 것은 6장에서 본 행렬 지수 $e^A$이다. 이 장에서는 먼저 행렬 함수를 정의하는 몇 가지 동치인 방법과 삼각 행렬의 함수가 갖는 구조를 보고, 행렬 지수를 계산하는 표준 방법인 거듭제곱 조정과 제곱의 원리를 Padé 근사와 함께 자세히 다룬다. 차수와 제곱 횟수를 고르는 데 노름이 필요하므로, 행렬을 만들지 않고 1-노름을 추정하는 Hager와 Higham의 방법도 함께 본다. 마지막으로 4장에서 소개한 Sylvester 방정식의 수치적 해법을 블록 알고리즘의 관점에서 다시 살핀다.}

\section{Definitions of Matrix Functions}
\label{sec:matrix-function-def}

멱급수 $f(z)=\sum c_kz^k$의 수렴 반지름이 $A$의 스펙트럼 반지름보다 크면 $f(A)=\sum c_kA^k$로 정의할 수 있다. $e^A=\sum A^k/k!$이 그 예이다. 더 일반적으로, $A=XJX^{-1}$을 Jordan 형식이라 할 때 크기 $m$인 Jordan 블록 $J_\lambda$에 대해
\[
f(J_\lambda)=\begin{pmatrix}f(\lambda)&f'(\lambda)&\cdots&f^{(m-1)}(\lambda)/(m-1)!\\&f(\lambda)&\ddots&\vdots\\&&\ddots&f'(\lambda)\\&&&f(\lambda)\end{pmatrix}
\]
로 두고 $f(A)=Xf(J)X^{-1}$로 정의할 수도 있다. 5장의 Cauchy 적분 공식을 행렬로 옮긴 $f(A)=(2\pi i)^{-1}\oint f(z)(zI-A)^{-1}dz$도 같은 값을 준다. 이 정의들은 모두 일치하며, $f(A)$는 $A$와 교환하는 행렬이고, $A$의 고유값 $\lambda$에 대해 $f(\lambda)$가 $f(A)$의 고유값이다\cite{higham2008functions}.

정의는 계산법을 바로 주지 않는다. $A$가 대각화 가능하더라도 $X$의 조건수가 크면 $Xf(\Lambda)X^{-1}$은 큰 오차를 만들고, Jordan 형식은 반올림 앞에서 불안정해 계산할 수 없다. 대칭 행렬이라면 직교 대각화 $A=Q\Lambda Q^{\mathsf T}$로 $f(A)=Qf(\Lambda)Q^{\mathsf T}$를 안정하게 계산할 수 있지만, 일반 행렬에서는 다른 방법이 필요하다.

\subsection{Functions of Triangular Matrices}

4장의 Schur 분해 $A=QTQ^{\mathsf H}$를 쓰면 $f(A)=Qf(T)Q^{\mathsf H}$이고, 문제는 위삼각 행렬의 함수 $F=f(T)$로 줄어든다. $F$도 위삼각이고 대각은 $f(t_{ii})$이다. 대각 위의 원소는 $FT=TF$라는 교환 관계에서 하나씩 정해지는데, $2\times2$의 경우를 보면
\[
f\begin{pmatrix}\lambda_1&t\\0&\lambda_2\end{pmatrix}=\begin{pmatrix}f(\lambda_1)&t\,\dfrac{f(\lambda_2)-f(\lambda_1)}{\lambda_2-\lambda_1}\\0&f(\lambda_2)\end{pmatrix}
\]
이다. 대각 위의 원소에 \term{나눗셈 차분}{divided difference}이 나타난다. 두 고유값이 가까우면 이 차분은 9장의 소거를 일으킨다. $f=\exp$, $\lambda_1=1$, $\lambda_2=1+10^{-10}$이면 공식대로 계산한 차분의 상대오차가 $4.9\times10^{-7}$이지만, 지수함수의 성질을 써서 $e^{\lambda_1}\,\operatorname{expm1}(\lambda_2-\lambda_1)/(\lambda_2-\lambda_1)$로 계산하면 $5.3\times10^{-17}$이다. 고유값이 모여 있으면 이런 소거가 행렬 전체에서 일어나므로, Schur 형식 위의 계산은 고유값을 무리 지어 블록으로 다루는 정교한 방법을 필요로 한다.

\section{The Matrix Exponential}
\label{sec:expm}

\subsection{Difficulties}

Moler와 Van Loan의 유명한 논문의 제목처럼 행렬 지수를 계산하는 방법은 적어도 열아홉 가지가 있고, 그 대부분은 어떤 행렬에서는 실패한다\cite{moler2003nineteen}. Taylor 급수를 그대로 더하면 노름이 큰 행렬에서 3장의 $e^x$에서 본 소거가 일어나고, 고유값 분해는 고유벡터가 거의 평행하면 실패한다. 특히 \term{비정규}{non-normal} 행렬, 곧 $AA^{\mathsf H}\ne A^{\mathsf H}A$인 행렬에서는 고유값이 모두 음수여도 $\|e^{tA}\|$가 $t$에 대해 한동안 크게 자랐다가 줄어드는 \term{혹}{hump}이 나타날 수 있고, 이것이 계산을 어렵게 만든다\cite{vanloan1977}.

\subsection{Scaling and Squaring}

가장 신뢰받는 방법은 $e^A=\bigl(e^{2^{-s}A}\bigr)^{2^s}$라는 항등식을 쓴다. 먼저 $A$를 $2^s$로 나누어 노름을 작게 만든 뒤 $e^{2^{-s}A}$를 근사하고 그 결과를 $s$번 제곱한다. 이것을 \term{거듭제곱 조정과 제곱}{scaling and squaring}이라 한다. 노름이 작은 행렬의 지수는 유리함수로 근사하는데, $e^x$의 $[m/m]$ \term{Padé 근사}{Padé approximant} $r_m(x)=p_m(x)/p_m(-x)$가 표준이다. Padé 근사는 같은 차수의 Taylor 다항식보다 원점 근처에서 훨씬 정확하고(30장에서 다룬다), $p_m(x)=\sum_jb_jx^j$에서 홀수 차 항과 짝수 차 항을 나누어 $U$, $V$라 하면 $r_m(B)=(V-U)^{-1}(V+U)$이다. 행렬 곱 몇 번과 연립방정식 풀기 하나로 계산된다.

Higham은 차수 $m$과 제곱 횟수 $s$를 고르는 기준을 후진 오차로 세웠다. $r_m(B)=e^{B+\Delta B}$로 쓸 때 $\|\Delta B\|/\|B\|\le u$가 되게 하는 $\|B\|$의 상한 $\theta_m$을 미리 계산해 두고, $\|2^{-s}A\|\le\theta_m$이 되는 가장 싼 $(m,s)$를 고르는 것이다\cite{higham2005expm}. 그러면 Padé 근사의 절단 오차만 보았을 때 계산한 근사는 $A$를 상대적으로 $u$ 이하만큼 바꾼 행렬의 정확한 지수이다. 근사를 평가하고 제곱하는 동안의 반올림 오차는 이 보장에 들어 있지 않다. 제곱 횟수가 너무 많으면 제곱하면서 반올림 오차가 커지므로, 제곱을 줄이는 것도 정확도에 중요하다. Al-Mohy와 Higham은 노름 $\|A\|$ 대신 $\|A^k\|^{1/k}$를 써서 판단을 정교하게 했다\cite{almohy2009expm}. 비정규 행렬에서는 $\|A^k\|^{1/k}$가 $\|A\|$보다 훨씬 작을 수 있어서 조정과 제곱을 덜 하고 오차도 줄어든다.

\begin{instardust}[\code{csrc/linalg/ops/expm.cpp}]
LAPACK에는 행렬 지수가 없으므로 stardust는 Al-Mohy와 Higham의 알고리즘을 직접 구현했다. scipy와 같은 알고리즘이다. 차수는 $3,5,7,9,13$ 가운데서 고르고, $\theta_3\approx0.0150$, $\theta_5\approx0.254$, $\theta_7\approx0.950$, $\theta_9\approx2.10$, $\theta_{13}=4.25$를 쓴다. 결과는 $(V-U)^{-1}(V+U)$ 대신 수학적으로 같은 $I+2(V-U)^{-1}U$로 만드는데, 단위 행렬을 작은 보정에 정확히 더하므로 노름이 작을 때 대각의 마지막 비트를 잃지 않는다. 설계 문서에 따르면 처음 꼴로는 $\|A\|\approx10^{-3}$에서 마지막 비트를 잃었다. 삼각 행렬이면 제곱할 때마다 대각을 정확한 $e^{2^{-k}a_{ii}}$로, 첫 부대각을 $2\times2$ 블록의 정확한 식으로 다시 쓴다. 그 식은 두 고유값이 가까우면 소거가 없는 꼴 $t\,e^{(\lambda_1+\lambda_2)/2}\operatorname{sinhc}((\lambda_2-\lambda_1)/2)$를, 멀면 나눗셈 차분을 쓴다. 가까울 때 나눗셈 차분을 쓰는 scipy는 $\begin{psmallmatrix}10&1\\0&10+10^{-12}\end{psmallmatrix}$의 $(0,1)$ 원소에서 상대오차 $2\times10^{-5}$를 낸다고 설계 문서는 적는다.
\end{instardust}

\section{1-Norm Estimation}
\label{sec:norm-estimation}

$\|A^k\|_1$을 알려면 $A^k$를 만들어야 하는데, 그것은 행렬 곱이어서 비싸다. 다행히 1-노름은 행렬 없이, 행렬--벡터 곱 몇 번으로 좋은 추정을 얻을 수 있다. $\|B\|_1=\max_{\|x\|_1\le1}\|Bx\|_1$이고, 최대는 단위 벡터 $e_j$ 가운데 하나에서 얻어진다. Hager는 이 최대화 문제에 볼록 함수의 경사를 따라 꼭짓점을 옮겨 가는 방법을 적용했다\cite{hager1984}. $x$에서 $y=Bx$를 계산하고 $\xi=\operatorname{sign}(y)$, $z=B^{\mathsf T}\xi$를 구하면, $z$의 가장 큰 성분의 자리 $j$가 $\|B\,\cdot\|_1$을 더 크게 하는 방향이므로 $x=e_j$로 옮긴다. 더 이상 커지지 않으면 멈춘다. 대개 몇 번 만에 수렴하고, 결과는 참값의 하한이지만 실제로는 참값과 같거나 매우 가까운 경우가 대부분이다.

\begin{algorithm}[t]
\begin{algorithmic}[1]
\Require 곱 $x\mapsto Bx$와 $y\mapsto B^{\mathsf T}y$를 계산하는 방법
\Ensure $\|B\|_1$의 하한 $\gamma$
\State $x\gets(1/n,\dots,1/n)^{\mathsf T}$
\Repeat
  \State $y\gets Bx$, $\gamma\gets\|y\|_1$
  \State $\xi\gets\operatorname{sign}(y)$, $z\gets B^{\mathsf T}\xi$
  \State $j\gets\arg\max_i|z_i|$
  \If{$|z_j|\le z^{\mathsf T}x$} \Return $\gamma$ \Comment{더 이상 커지지 않는다}
  \EndIf
  \State $x\gets e_j$
\Until{정해진 횟수}
\State \Return $\gamma$
\end{algorithmic}
\caption{Hager의 1-노름 추정}
\label{alg:hager}
\end{algorithm}

Higham은 이 방법이 실패하는 예를 막는 보조 추정을 덧붙여 LAPACK의 \code{lacon}과 그 후속 \code{lacn2}로 만들었고\cite{higham1988lacon}, Higham과 Tisseur는 여러 벡터를 한꺼번에 다루는 블록 판을 만들었다\cite{higham2000block}. $B=A^{-1}$이면 곱 $A^{-1}x$는 19장의 LU 분해로 $O(n^2)$에 계산되므로, $\|A^{-1}\|_1$을 역행렬 없이 추정할 수 있다. LAPACK의 조건수 추정 \code{gecon}이 이것이다.

\begin{instardust}[\code{csrc/linalg/ops/expm.cpp}]
stardust의 \code{expm}은 $\|A^2\|$부터 $\|A^6\|$까지는 거듭제곱을 실제로 만들어 정확한 노름을 쓰고, 이미 아는 노름의 곱으로 만든 상한이 판정을 통과하면 그 거듭제곱조차 만들지 않는다. $n\ge100$이면 $\|A^8\|$과 $\|A^{10}\|$은 \code{lacn2}로 추정해 행렬 곱 대신 이미 있는 $A^4$와 $A^6$에 행렬--벡터 곱을 두 번씩 한다. 설계 문서는 차수 $13$이 필요한 큰 행렬에서 이 두 곱이 호출 시간의 $12$--$14\%$였고, 추정으로 바꾸어도 잰 $60$개 행렬 모두에서 고른 제곱 횟수가 같았다고 적는다.
\end{instardust}

\section{The Sylvester Equation}
\label{sec:sylvester}

4장에서 본 Sylvester 방정식 $AX+XB=C$는 행렬 함수와 깊이 연결되어 있다. 예컨대 블록 위삼각 행렬의 함수의 대각 밖 블록은 Sylvester 방정식의 해로 주어지고, 행렬 제곱근의 Fréchet 미분 $L$은 $XL+LX=E$를 푼다. 4장의 Bartels--Stewart 방법은 $A$와 $B$를 Schur 형식으로 바꾸어 방정식을 삼각 꼴로 만든 뒤 원소를 하나씩 풀었다\cite{bartels1972stewart}. 그 마지막 단계는 원소마다 앞서 구한 원소들의 일차결합을 빼는 2층 연산의 연속이어서 큰 행렬에서는 메모리에 묶인다. 19장의 블록 LU처럼 해를 블록으로 나누고 블록 사이의 갱신을 행렬 곱으로 하면 3층 BLAS의 속도를 얻는데, Jonsson과 Kågström은 이것을 재귀적으로 블록을 나누는 꼴로 정리했다\cite{jonsson2002recsy}. LAPACK의 \code{trsyl3}은 넘침을 막는 척도 조정을 갖춘 비재귀 블록 판이다.

\begin{instardust}[\code{docs/design/linalg.md} §1.2]
stardust의 \code{solve\_sylvester}는 두 행렬을 \code{gees}로 Schur 형식으로 바꾼 뒤 LAPACK의 블록 삼각 Sylvester 풀이 \code{trsyl3}을 부르고 결과를 Schur 벡터로 되돌린다. 행렬 곱은 \code{sd.matmul}을 거치지 않고 BLAS의 \code{gemm}을 직접 부르는데, 작업 버퍼가 이미 C++ 안에 F 순서로 있기 때문이다.
\end{instardust}

\section{Concluding Remarks}

행렬 함수는 스칼라 함수의 단순한 확장처럼 보이지만, 그 계산은 비정규성과 가까운 고유값이라는 행렬만의 어려움을 안고 있다. 거듭제곱 조정과 제곱은 이 어려움을 후진 오차의 관점에서 정면으로 다루어 차수와 제곱 횟수를 절단 오차가 반올림 수준의 후진 오차가 되도록 고른다. 그 판단에 필요한 노름은 행렬 곱 대신 몇 번의 행렬--벡터 곱으로 추정한다. 다음 장에서는 원소 대부분이 $0$인 희소 행렬로 넘어가 그 구조를 살리는 저장 방식과 분해를 다룬다.

행렬 함수의 이론과 계산은 Higham의 책\cite{higham2008functions}이 표준이다.

\chapterreferences
